% Transporting the Group Law (Why It Is Forced)  (merge of source §12 + §19:
% the transported law, the proof that it equals chord-tangent via Thm 10.5, the
% "no arbitrary convention" logic, and intrinsicness/uniqueness)
\chapter{Transporting the Group Law (Why It Is Forced)}
\label{ch:transport}

Everything is now in place to prove the main theorem. The isomorphism \(\Phi\)
moves the torus group law onto the cubic; the collinearity criterion of
\cref{ch:lines_divisors} shows that the moved law is the chord--tangent
construction. This chapter assembles the two, isolates the logical mechanism that
makes the chord--tangent rule unavoidable, and observes that the resulting group
law is intrinsic to the pair \((E, O)\) --- independent of which lattice and
embedding produced it.

\section{The Transported Group Law}
\label{sec:tr-transported}

The torus already carries an abelian group law, \([u] + [v] = [u+v]\), with no
freedom: it is addition in \(\C\) (\cref{prop:tg-welldef}). The biholomorphism
\(\Phi : T \xrightarrow{\sim} E\) of \cref{thm:emb-iso} carries it to \(E\) by the
only rule that makes \(\Phi\) a group isomorphism.

\begin{definition}
\label{def:tr-law}
The \emph{transported group law} on \(E\) is
\[
  P \oplus Q = \Phi\bigl( \Phi^{-1}(P) + \Phi^{-1}(Q) \bigr),
\]
with identity \(O = \Phi(0) = [0:1:0]\) and inverse \(\ominus P = \Phi(-\Phi^{-1}(P))\).
In affine coordinates, by \cref{sec:emb-reflection}, the inverse is reflection,
\(\ominus(x, y) = (x, -y)\).
\end{definition}

By construction \((E, \oplus)\) is an abelian group and \(\Phi\) is an
isomorphism of groups. The substantive question is geometric:

\begin{quote}
What is \(P \oplus Q\) in terms of the cubic alone?
\end{quote}

\section{The Chord--Tangent Construction}
\label{sec:tr-chord-tangent}

\begin{definition}
\label{def:tr-chordtangent}
The \emph{chord--tangent operation} \(\boxplus\) on \(E\) is defined as follows.
For \(P, Q \in E\), let \(L\) be the line through \(P\) and \(Q\) --- the tangent to
\(E\) at \(P\) if \(P = Q\) --- and let \(R\) be the third point of \(L \cap E\)
(\cref{prop:lin-bezout}). Then
\[
  P \boxplus Q = \ominus R,
\]
the reflection of \(R\). Here \(\ominus R\) is the third intersection of the line
\(RO\) with \(E\), which in affine coordinates is \((x_R, -y_R)\).
\end{definition}

This is the recipe of \cref{sec:intro-construction}: draw the chord (or tangent),
take the third intersection, reflect. The main theorem is that it agrees with the
transported law.

\begin{theorem}[Chord--tangent is the group law]
\label{thm:tr-main}
Under the Weierstrass embedding \(\Phi : \C/\Lambda \to E\),
\(z \mapsto [\wp(z):\wp'(z):1]\), the transported group law equals the
chord--tangent operation:
\[
  P \oplus Q = P \boxplus Q \qquad\text{for all } P, Q \in E .
\]
Consequently the chord--tangent construction is an abelian group law on \(E\) with
identity \(O\), and \(\Phi\) is an isomorphism of groups
\((\C/\Lambda, +) \cong (E, \boxplus)\). This is the group-theoretic assertion of
\cref{thm:main}.
\end{theorem}

\begin{proof}
Let \(L\) be the line through \(P\) and \(Q\) (tangent if \(P = Q\)), meeting \(E\)
in \(P + Q + R\) with multiplicity. By the collinearity criterion
\cref{thm:lin-collinear}, the three points sum to the identity in the transported
law:
\[
  P \oplus Q \oplus R = O,
  \qquad\text{hence}\qquad
  P \oplus Q = \ominus R .
\]
The group inverse \(\ominus R\) is the reflection of \(R\)
(\cref{def:tr-law,sec:emb-reflection}), which is exactly \(P \boxplus R\)'s
defining reflection, i.e.\ \(P \oplus Q = \ominus R = P \boxplus Q\). The identity,
inverse, and tangent cases are the same statement specialized: when \(P = Q\) the
line is the tangent and \(L \cdot E = 2P + R\), giving \(2P = \ominus R\); when
\(Q = \ominus P\) the line \(PQ\) is vertical through \(O\) and \(R = O\), giving
\(P \oplus (\ominus P) = O\); when \(Q = O\) the line \(PO\) meets \(E\) at
\(P, \ominus P, O\), giving \(P \oplus O = P\). All are instances of
\cref{thm:lin-collinear}, worked out one by one in \cref{ch:chord_tangent_cases}.
\end{proof}

\section{Why the Rule Is Forced}
\label{sec:tr-forced}

The proof exposes the exact mechanism. The chord--tangent law is not a clever
convention; it is the only possible answer, dictated by three inputs, none of
which involves a choice.

\begin{enumerate}
\item \textbf{The torus has a group law.} \([u]+[v]=[u+v]\) is addition in \(\C\),
  forced (\cref{prop:tg-welldef}).
\item \textbf{The embedding transfers it.} Requiring \(\Phi\) to be a group
  isomorphism leaves no freedom in \(\oplus\) (\cref{def:tr-law}).
\item \textbf{Lines see that law.} A line cuts \(E\) in a divisor \(P+Q+R-3O\)
  that is principal (\cref{thm:lin-principal}), and Abel's law makes its points
  sum to \(O\) (\cref{thm:lin-collinear}). So a line through \(P, Q\) determines
  \(R = \ominus(P \oplus Q)\), and the sum is recovered by the one further forced
  step --- negate, i.e.\ reflect.
\end{enumerate}

\begin{keyidea}{No arbitrary convention}
The reflection in the chord--tangent rule is often presented as a definition to be
accepted. It is not optional. A line gives the symmetric relation
\(P + Q + R = O\), because line divisors are principal and principal divisors have
Abel sum zero. Solving this for \(P \oplus Q\) forces
\[
  P \oplus Q = \ominus R,
\]
and \(\ominus\) is reflection because the torus inverse \(z \mapsto -z\) is
\((x,y) \mapsto (x,-y)\). Chord, tangent, third point, reflect: each step is
compelled.
\end{keyidea}

\section{The Group Law Is Intrinsic}
\label{sec:tr-intrinsic}

A subtle point deserves emphasis. The construction passed through a lattice
\(\Lambda\) and an embedding \(\Phi\), but the resulting law does not depend on
them.

\begin{proposition}
\label{prop:tr-intrinsic}
The chord--tangent law \(\boxplus\) depends only on the cubic \(E\) and the choice
of identity \(O\); it is independent of the lattice \(\Lambda\) and the embedding
\(\Phi\) used to construct it.
\end{proposition}

\begin{proof}
\Cref{def:tr-chordtangent} refers only to lines in \(\CP^2\), their
intersections with \(E\), and the point \(O\) --- all intrinsic to \((E, O)\). No
lattice or \(\wp\) appears. \Cref{thm:tr-main} shows this intrinsic operation
coincides with any transported law \(\oplus\) built from a \(\Phi\); hence all such
transported laws agree with \(\boxplus\) and therefore with each other.
\end{proof}

In other words, the abelian group is attached to the geometric object \((E, O)\)
itself. The uniformization \(\C/\Lambda \cong E\) explains \emph{why} the law is a
group and why it is commutative and associative --- properties that are opaque from
the cubic but transparent on the torus --- yet the law it computes is the one a
geometer reading only the plane curve would write down.

\begin{remark}
\label{rem:tr-associativity}
Associativity of \(\boxplus\) is the property hardest to see directly on the cubic;
the classical proofs invoke the Cayley--Bacharach theorem. Here it is free: \(\oplus\)
is associative because addition in \(\C\) is, and \(\boxplus = \oplus\). The torus
is, in this sense, the explanation of associativity.
\end{remark}

\section{Summary}
\label{sec:tr-summary}

Transporting the torus law by \(\Phi\) gives a group \((E, \oplus)\) with identity
\(O\) and reflection as inverse (\cref{def:tr-law}). By the collinearity criterion,
this transported law is exactly the chord--tangent construction
(\cref{thm:tr-main}): a line through \(P, Q\) meets \(E\) again at
\(R = \ominus(P\oplus Q)\), so \(P \oplus Q\) is the reflection of \(R\). The rule
is forced by three choiceless facts --- the torus group law, its transfer by
\(\Phi\), and the principality of line divisors (\cref{sec:tr-forced}) --- and the
resulting law is intrinsic to \((E, O)\), with associativity inherited from \(\C\)
(\cref{prop:tr-intrinsic,rem:tr-associativity}). The remaining chapters make the
construction explicit: the pictorial dictionary and the case analysis
(\cref{ch:chord_tangent_cases}), then closed-form coordinates
(\cref{ch:coordinate_formulas}).

\section*{Exercises}
\addcontentsline{toc}{section}{Exercises}

\begin{exercise}
\label{ex:tr-unique-transport}
Show that \(\oplus\) in \cref{def:tr-law} is the unique binary operation on \(E\)
making \(\Phi\) a group homomorphism, and verify directly that it is abelian with
identity \(O\) and inverse \(\ominus P = \Phi(-\Phi^{-1}(P))\).
\end{exercise}

\begin{exercise}
\label{ex:tr-identity-case}
Work out the identity case from \cref{thm:tr-main}: show the line through \(P\) and
\(O\) meets \(E\) at \(P\), \(\ominus P\), and \(O\), and deduce \(P \oplus O = P\)
directly from \cref{thm:lin-collinear}.
\end{exercise}

\begin{exercise}
\label{ex:tr-associativity-free}
Explain in one paragraph why associativity of the chord--tangent law is immediate
from \cref{thm:tr-main}, and state precisely which property of \(\C\) it reduces
to. Why is the direct geometric proof hard?
\end{exercise}

\begin{exercise}
\label{ex:tr-intrinsic}
Suppose two lattices \(\Lambda, \Lambda'\) and embeddings \(\Phi, \Phi'\) both
present the \emph{same} cubic \(E\) (with the same \(O\)). Using
\cref{prop:tr-intrinsic}, show the two transported laws coincide. What does this
say about the relationship between \(\Lambda\) and \(\Lambda'\) (cf.\
\cref{prop:lat-normal-form})?
\end{exercise}

\begin{exercise}
\label{ex:tr-translation}
Show that for fixed \(A \in E\), the translation \(P \mapsto P \oplus A\) is a
biholomorphism of \(E\) (an automorphism of the Riemann surface), corresponding
under \(\Phi\) to \(z \mapsto z + \Phi^{-1}(A)\) on \(T\)
(\cref{prop:tg-lie}). Deduce that \(E\) acts on itself by holomorphic
translations, so \(E\) is a complex Lie group.
\end{exercise}
